\documentclass[11pt,a4paper]{article}

\setlength{\hoffset}{1in}
\setlength{\oddsidemargin}{-1in}
\setlength{\voffset}{0.9in}
\setlength{\topmargin}{-1.6in}
\setlength{\textwidth}{15.92cm}
\setlength{\textheight}{25cm}
\linespread{1.1}

\usepackage[utf8]{inputenc}
\usepackage{amsmath}
\usepackage{amsfonts}
\usepackage{amssymb}
\usepackage{makeidx}
\usepackage[all]{xy}


\begin{document}

\begin{center}
SHENZHEN UNIVERSITY \\
INSTITUTE FOR ADVANCED STUDY \\
Differential Equations, Autumn 2026
Assignment 2, due on Thursday, October 12, 2026, at 6:00 PM,

\end{center}


\bigskip
\noindent
{\bf Problem 1.}  (Problems 7,10,11 of Exercises 2.2) Solve the following equations:

1) $x\dfrac{dy}{dx}=\dfrac{1}{y^3}$.
\medskip
\par\noindent
\textbf{Solution.}\par
\[
y^3\,dy=\frac{1}{x}\,dx
\]
\[
\int y^3\,dy=\int\frac{1}{x}\,dx
\]
\[
\frac{1}{4}y^4=\ln|x|+C
\]
\[
\ln|x|-\frac{1}{4}y^4=C
\]
\medskip
2) $\dfrac{dy}{dx}=\dfrac{x}{y^2\sqrt{1+x}}$.
\medskip
\par\noindent
\textbf{Solution.}\par
\[
y^2\,dy=
\left(\frac{1+x}{\sqrt{1+x}}-\frac{1}{\sqrt{1+x}}\right)dx
\]
\[
\int y^2\,dy=
\int\left(\frac{1+x}{\sqrt{1+x}}-\frac{1}{\sqrt{1+x}}\right)d(x+1)
\]
\[
\frac{1}{3}y^3=
\frac{2}{3}(x+1)^{3/2}-2(x+1)^{1/2}+C
\]
\[
y^3-2(x+1)^{3/2}+6(x+1)^{1/2}=C_1
\]
\medskip
3) $x\dfrac{dv}{dx}=\dfrac{1-4v^2}{3v}$.
\medskip
\par\noindent
\textbf{Solution.}\par
\[
\frac{3v}{1-4v^2}\,dv=\frac{1}{x}\,dx
\]
\[
\int\frac{3v}{1-4v^2}\,dv=\int\frac{1}{x}\,dx
\]
\[
\int\frac{3}{4}
\left(\frac{1}{1-2v}-\frac{1}{1+2v}\right)dv
=\int\frac{1}{x}\,dx
\]
\[
\frac{3}{4}\left[
-\frac{1}{2}\int\frac{1}{1-2v}\,d(1-2v)
-\frac{1}{2}\int\frac{1}{1+2v}\,d(1+2v)
\right]
=\int\frac{1}{x}\,dx
\]
\[
-\frac{3}{8}\left[\ln|1-2v|+\ln|1+2v|\right]
=\ln|x|+C
\]
\[
-\frac{3}{8}\ln|1-4v^2|=\ln|x|+C
\]
\[
\ln|1-4v^2|=\ln\left(|x|^{-8/3}\right)+C
\]
\[
|1-4v^2|=C_2|x|^{-8/3}
\]
\[
1-4v^2=C_2x^{-8/3}
\]
\medskip
\bigskip
\noindent
{\bf Problem 2.}  (Problems 20 and 21 of Exercises 2.2)
Solve the following initial value problems:

1) $x^2\dfrac{dy}{dx}=\dfrac{4x^2- x-2}{(x+1)(y+1)}, \quad y(1)=1$.
\medskip
\par\noindent
\textbf{Solution.}\par
\[
x^2\frac{dy}{dx}
=\frac{(x+1)(4x-5)+3}{(x+1)(y+1)}
\]
\[
x^2\frac{dy}{dx}
=\frac{1}{y+1}\left(4x-5+\frac{3}{x+1}\right)
\]
\[
(y+1)\,dy
=\left[\frac{4}{x}-\frac{5}{x^2}
+\frac{3}{x^2(x+1)}\right]dx
\]
\[
\int(y+1)\,dy
=\int\left[\frac{4}{x}-\frac{5}{x^2}
-\frac{3}{x}+\frac{3}{x^2}+\frac{3}{x+1}\right]dx
\]
\[
\int(y+1)\,dy
=\int\left(\frac{1}{x}-\frac{2}{x^2}
+\frac{3}{x+1}\right)dx
\]
\[
\frac{1}{2}y^2+y
=\ln|x|+\frac{2}{x}+3\ln|x+1|+C
\]
Substitute $y(1)=1$
\[
\frac{1}{2}+1=2+3\ln 2+C
\]
\[
C=-\frac{1}{2}-\ln 8
\]
The solution is:
\[
\frac{1}{2}y^2+y
=\ln|x|+\frac{2}{x}+3\ln|x+1|-\frac{1}{2}-\ln 8
\]
\medskip
2) $\dfrac{1}{\theta}\dfrac{dy}{d\theta}=\dfrac{y \sin \theta}{y^2+1},   \quad y(\pi)=1$.
\medskip
\par\noindent
\textbf{Solution.}\par
\[
\left(y+\frac{1}{y}\right)dy=\theta\sin\theta\,d\theta
\]
\[
\int\left(y+\frac{1}{y}\right)dy
=\int\theta\sin\theta\,d\theta
\]
Suppose
\[
\begin{aligned}
u  &= \theta,  & \qquad dv &= \sin\theta\,d\theta,\\
du &= d\theta, & v         &= -\cos\theta.
\end{aligned}
\]
\[
\begin{aligned}
\int\theta\sin\theta\,d\theta
&=uv-\int v\,du\\
&=-\theta\cos\theta+\int\cos\theta\,d\theta\\
&=-\theta\cos\theta+\sin\theta+C
\end{aligned}
\]
\[
\ln|y|+\frac{1}{2}y^2
=-\theta\cos\theta+\sin\theta+C
\]
Substitute $y(\pi)=1$
\[
\frac{1}{2}=-\pi\cdot(-1)+C
\]
\[
C=\frac{1}{2}-\pi
\]
The solution is:
\[
\ln y+\frac{1}{2}y^2
=-\theta\cos\theta+\sin\theta+\frac{1}{2}-\pi
\]
\medskip
  

\bigskip
\noindent
{\bf Problem 3.}  (Problems 7 and 12 of Exercises 2.3)
Find the general solutions to the following equations:

1) $\dfrac{dy}{dx}-y-e^{3x}=0$.
\medskip
\par\noindent
\textbf{Solution.}\par
\[
\frac{dy}{dx}-y=e^{3x}
\]
\[
P(x)=-1,\qquad Q(x)=e^{3x}
\]
\[
e^{\int P(x)\,dx}=e^{\int(-1)\,dx}=e^{-x}
\]
\[
e^{-x}\frac{dy}{dx}-e^{-x}y=e^{2x}
\]
\[
\frac{d}{dx}\left(e^{-x}y\right)=e^{2x}
\]
\[
\int\frac{d}{dx}\left(e^{-x}y\right)dx
=\int e^{2x}\,dx
\]
\[
e^{-x}y=\frac{1}{2}e^{2x}+C
\]
\[
y=e^x\left(\frac{1}{2}e^{2x}+C\right)
\]
\[
y=\frac{1}{2}e^{3x}+Ce^x
\]
\medskip
2)  $\dfrac{dy}{dx}=x^2e^{-4x}-4y$.
\medskip
\par\noindent
\textbf{Solution.}\par
\[
\frac{dy}{dx}+4y=x^2e^{-4x}
\]
\[
P(x)=4,\qquad Q(x)=x^2e^{-4x}
\]
\[
e^{\int P(x)\,dx}=e^{4x}
\]
\[
e^{4x}\frac{dy}{dx}+4ye^{4x}=x^2
\]
\[
\frac{d}{dx}\left(e^{4x}y\right)=x^2
\]
\[
\int\frac{d}{dx}\left(e^{4x}y\right)dx
=\int x^2\,dx
\]
\[
e^{4x}y=\frac{1}{3}x^3+C
\]
\[
y=\frac{1}{e^{4x}}\left(\frac{1}{3}x^3+C\right)
\]
\medskip



\bigskip
\noindent
{\bf Problem 4.}
(Problem 25 of Exercises 2.3)  \\
\noindent
  Using definite integration, show that the solution to the initial value problem
 $$
 \dfrac{dy}{dx} + 2xy =1,  \quad y(2)=1,
 $$
 can be expressed as
 $$
 y(x) = e^{-x^2} \left( e^4 + \int_2^x e^{t^2} dt \right) .
 $$
\medskip
\par\noindent
\textbf{Solution.}\par
\[
P(x)=2x,\qquad Q(x)=1
\]
\[
e^{\int P(x)\,dx}=e^{\int 2x\,dx}=e^{x^2}
\]
\[
e^{x^2}\frac{dy}{dx}+e^{x^2}\cdot 2xy=e^{x^2}
\]
\[
\frac{d}{dx}\left(e^{x^2}y\right)=e^{x^2}
\]
\[
\int_2^x\frac{d}{dt}\left(e^{t^2}y(t)\right)dt
=\int_2^x e^{t^2}\,dt
\]
\[
e^{x^2}y(x)-e^{2^2}y(2)=\int_2^x e^{t^2}\,dt
\]
\[
e^{x^2}y(x)-e^4=\int_2^x e^{t^2}\,dt
\]
\[
y(x)=e^{-x^2}\left[e^4+\int_2^x e^{t^2}\,dt\right]
\]
\medskip




\bigskip
\noindent
{\bf Problem 5.}  (Problems 10,18, 20 of Exercises 2.4)
Determine whether the following equations are exact. If they are, then solve them.

1) $(2x+y)dx+(x-2y)dy=0$.
\medskip
\par\noindent
\textbf{Solution.}\par
\[
M(x,y)=2x+y,\qquad N(x,y)=x-2y
\]
\[
\frac{\partial M}{\partial y}=1,\qquad
\frac{\partial N}{\partial x}=1
\]
$\dfrac{\partial M}{\partial y}=\dfrac{\partial N}{\partial x}$,
while $M,N$ have continuous first partial derivatives.

Thus, it is an exact equation.
\[
\begin{aligned}
F&=\int M\,dx+g(y)\\
 &=\int(2x+y)\,dx+g(y)\\
 &=x^2+xy+g(y)
\end{aligned}
\]
\[
N=\frac{d}{dy}(x^2+xy)+g'(y)
\]
\[
\begin{aligned}
g(y)&=\int\left[N-\frac{d}{dy}(x^2+xy)\right]dy\\
    &=\int[x-2y-x]\,dy\\
    &=-y^2
\end{aligned}
\]
\[
F(x,y)=x^2+xy-y^2=C
\]
\medskip
2) $(2x+y^2-\cos (x+y))dx+(2xy-\cos (x+y)-e^y)dy=0$.
\medskip
\par\noindent
\textbf{Solution.}\par
\[
M(x,y)=2x+y^2-\cos(x+y)
\]
\[
N(x,y)=2xy-\cos(x+y)-e^y
\]
\[
\frac{\partial M}{\partial y}=2y+\sin(x+y),\qquad
\frac{\partial N}{\partial x}=2y+\sin(x+y)
\]
$\dfrac{\partial M}{\partial y}=\dfrac{\partial N}{\partial x}$,
while $M,N$ have continuous first derivatives.

Thus, it is an exact equation.
\[
\begin{aligned}
F&=\int M\,dx+g(y)\\
 &=\int[2x+y^2-\cos(x+y)]\,dx+g(y)\\
 &=x^2+xy^2-\sin(x+y)+g(y)
\end{aligned}
\]
\[
\begin{aligned}
N&=\frac{d}{dy}[x^2+xy^2-\sin(x+y)]+g'(y)\\
 &=2xy-\cos(x+y)+g'(y)
\end{aligned}
\]
\[
\begin{aligned}
g(y)&=\int[N-2xy+\cos(x+y)]\,dy\\
    &=\int[2xy-\cos(x+y)-e^y-2xy+\cos(x+y)]\,dy\\
    &=\int -e^y\,dy\\
    &=-e^y
\end{aligned}
\]
\[
F(x,y)=x^2+xy^2-\sin(x+y)-e^y=C
\]
\medskip
3)$\left(\dfrac{2}{\sqrt{1-x^2}}+y\cos (xy)\right)dx+\left(x\cos (xy)-y^{-1/3}\right)dy=0$.
\medskip
\par\noindent
\textbf{Solution.}\par
\[
M(x,y)=\frac{2}{\sqrt{1-x^2}}+y\cos(xy),\qquad
N(x,y)=x\cos(xy)-y^{-1/3}
\]
\[
\frac{\partial M}{\partial y}=\cos(xy)-xy\sin(xy)
\]
\[
\frac{\partial N}{\partial x}=\cos(xy)-xy\sin(xy)
\]
$\dfrac{\partial M}{\partial y}=\dfrac{\partial N}{\partial x}$,
while $M,N$ have continuous first derivatives.

Thus, it is an exact equation.
\[
\begin{aligned}
F&=\int M\,dx+g(y)\\
 &=\int\left[\frac{2}{\sqrt{1-x^2}}+y\cos(xy)\right]dx+g(y)\\
 &=\sin(xy)+2\arcsin x+g(y)
\end{aligned}
\]
\[
\begin{aligned}
N&=\frac{d}{dy}[\sin(xy)+2\arcsin x]+g'(y)\\
 &=x\cos(xy)+g'(y)
\end{aligned}
\]
\[
\begin{aligned}
g(y)&=\int[N-x\cos(xy)]\,dy\\
    &=\int[x\cos(xy)-y^{-1/3}-x\cos(xy)]\,dy\\
    &=-\frac{3}{2}y^{2/3}
\end{aligned}
\]
\[
F(x,y)=\sin(xy)+2\arcsin x-\frac{3}{2}y^{2/3}=C
\]
\medskip


\bigskip
\noindent
{\bf Problem 6.}  (Problem 26 of Exercises 2.4)
Solve the following initial problem.


$(\tan y-2)dx+(x\sec ^2y+1/y)dy=0, \quad \quad y(0)=1$
\medskip
\par\noindent
\textbf{Solution.}\par
\[
M(x,y)=\tan y-2,\qquad
N(x,y)=x\sec^2 y+\frac{1}{y}
\]
\[
\frac{\partial M}{\partial y}=\sec^2 y,\qquad
\frac{\partial N}{\partial x}=\sec^2 y
\]
$\dfrac{\partial M}{\partial y}=\dfrac{\partial N}{\partial x}$,
while $M,N$ have continuous first derivatives.

Thus, it is an exact equation.
\[
\begin{aligned}
F&=\int M\,dx+g(y)\\
 &=\int(\tan y-2)\,dx+g(y)\\
 &=x\tan y-2x+g(y)
\end{aligned}
\]
\[
\begin{aligned}
N&=\frac{d}{dy}(x\tan y-2x)+g'(y)\\
 &=x\sec^2 y+g'(y)
\end{aligned}
\]
\[
\begin{aligned}
g(y)&=\int(N-x\sec^2 y)\,dy\\
    &=\int\left(x\sec^2 y+\frac{1}{y}-x\sec^2 y\right)dy\\
    &=\ln|y|
\end{aligned}
\]
\[
F(x,y)=x\tan y-2x+\ln|y|=C
\]
Substitute $y(0)=1$
\[
C=0
\]
\[
F(x,y)=x\tan y-2x+\ln y=0
\]
\medskip

\bigskip
\noindent
{\bf Problem 7.}
(Problem 30 of Exercises 2.4) Consider the equation
$$
(5x^2y+6x^3y^2+4xy^2)dx+(2x^3+3x^4y+3x^2y)dy=0
$$
(a) Show that the equation is not exact.  \\
\noindent
(b) Multiply the equation by $x^ny^m$ and determine values for $n$ and $m$ that make the resulting equation exact. \\
\noindent
(c) Use the solution of the resulting exact equation to solve the original equation.

\medskip
\par\noindent
\textbf{Solution.}\par
\noindent
(a)
\[
M(x,y)=5x^2y+6x^3y^2+4xy^2,\qquad
N(x,y)=2x^3+3x^4y+3x^2y
\]
\[
\frac{\partial M}{\partial y}=5x^2+12x^3y+8xy,\qquad
\frac{\partial N}{\partial x}=6x^2+12x^3y+6xy
\]
$\dfrac{\partial M}{\partial y}\ne\dfrac{\partial N}{\partial x}$,
it is not an exact equation.

\medskip
\noindent
(b)
\[
\begin{aligned}
M_1(x,y)&=x^ny^mM(x,y)\\
        &=5x^{2+n}y^{1+m}+6x^{3+n}y^{2+m}
          +4x^{1+n}y^{2+m}
\end{aligned}
\]
\[
\begin{aligned}
N_1(x,y)&=x^ny^mN(x,y)\\
        &=2x^{3+n}y^m+3x^{4+n}y^{1+m}
          +3x^{2+n}y^{1+m}
\end{aligned}
\]
\[
\begin{aligned}
\frac{\partial M_1}{\partial y}
&=5(1+m)x^{2+n}y^m+6(2+m)x^{3+n}y^{1+m}\\
&\quad+4(2+m)x^{1+n}y^{1+m}
\end{aligned}
\]
\[
\begin{aligned}
\frac{\partial N_1}{\partial x}
&=2(3+n)x^{2+n}y^m+3(4+n)x^{3+n}y^{1+m}\\
&\quad+3(2+n)x^{1+n}y^{1+m}
\end{aligned}
\]
\[
\begin{cases}
5(1+m)=2(3+n),\\
6(2+m)=3(4+n),\\
4(2+m)=3(2+n)
\end{cases}
\quad\Longrightarrow\quad
\begin{cases}
m=1,\\
n=2
\end{cases}
\]

\medskip
\noindent
(c)
\[
M_1(x,y)=5x^4y^2+6x^5y^3+4x^3y^3
\]
\[
N_1(x,y)=2x^5y+3x^6y^2+3x^4y^2
\]
\[
\begin{aligned}
F&=\int M_1\,dx+g(y)\\
 &=\int(5x^4y^2+6x^5y^3+4x^3y^3)\,dx+g(y)\\
 &=x^5y^2+x^6y^3+x^4y^3+g(y)
\end{aligned}
\]
\[
\begin{aligned}
N_1&=\frac{d}{dy}(x^5y^2+x^6y^3+x^4y^3)+g'(y)\\
   &=2x^5y+3x^6y^2+3x^4y^2+g'(y)
\end{aligned}
\]
\[
\begin{aligned}
g(y)&=\int[N_1-(2x^5y+3x^6y^2+3x^4y^2)]\,dy\\
    &=0
\end{aligned}
\]
\[
F(x,y)=x^5y^2+x^6y^3+x^4y^3=C
\]
\medskip



\end{document}












